\chapter{Quickstart}

\section{Getting a Poplog}

Three ways in, fastest first.

\textbf{1. A relocatable binary} (about a 2\,MB download). This unpacks a
ready-built \code{basepop11} plus libraries under
\code{\textasciitilde/.local/share/pop11-skill}, and finishes with a live
smoke test:

\begin{sh}
curl -fsSL https://raw.githubusercontent.com/IoTone/poplog/master/tools/install-skill.sh | sh
\end{sh}

\textbf{2. From source.} Poplog bootstraps from a small seed binary,
\code{corepop}, which is vendored for the ported platforms:

\begin{sh}
git clone https://github.com/IoTone/poplog && cd poplog
./configure && make all
\end{sh}

\code{make all} builds the core system, then compiles the Prolog, Common Lisp
and Standard ML front-ends and saves each as an image under
\code{target/psv/}. See \code{INSTALL} for the details.

\textbf{3. Nix.} A flake builds and bootstraps everything from source with no
manual seed or toolchain setup:

\begin{sh}
nix run .#pop11          # or .#prolog / .#clisp / .#pml
nix shell .#poplog       # put all five front-ends on $PATH
\end{sh}

\section{The first session}

\begin{repl}
$ ./poplog pop11
Sussex Poplog (Version 16.019997 ...)
Copyright (c) 1982-1999 University of Sussex.  All rights reserved.

Setpop
: 6 * 7 =>
** 42
\end{repl}

The prompt is \code{:} and the workhorse is \code{=>}, which prints whatever
the preceding expression left, prefixed with \code{**}. Two more you will use
constantly: \code{npr(x)} prints \code{x} and a newline, and \code{==>} is
\code{=>} with fuller structure printing.

Statements end with a semicolon. An expression followed by \code{=>} does not
need one --- \code{=>} terminates it.

When something goes wrong you get a \emph{mishap}: a diagnostic naming the
error, the culprits, and the chain of procedures that was running.

\begin{repl}
: 'three' + 4 =>
;;; MISHAP - NUMBER(S) NEEDED
;;; INVOLVING:  'three' 4
;;; PRINT DOING
;;; DOING    :  + pop_setpop_compiler
\end{repl}

The session survives. Fix the expression and carry on --- this is the normal
working rhythm, not an exceptional one.

\section{Running files}

A Pop-11 source file is just a sequence of statements, so there is no
\code{main}:

\begin{sh}
./poplog basepop11 myprog.p        # compile and run, then enter the REPL
./poplog basepop11 myprog.p </dev/null   # ...or run and exit
\end{sh}

Inside a session, \code{load 'myprog.p';} does the same thing, and
\code{uses foo;} loads library \code{foo} from the library search list if it
is not already loaded --- this is how you reach \code{lib json}, \code{lib
crypto}, \code{objectclass}, and the rest.

\section{Documentation, in the system}

Poplog ships its own manual --- over 900 files --- and it is queryable from
the prompt. The three kinds matter:

\begin{itemize}
\item \code{help \emph{name}} --- practical, task-shaped. Start here.
\item \code{ref \emph{name}} --- the exhaustive reference for a primitive or
  a subsystem.
\item \code{teach \emph{name}} --- tutorials, many of them the original
  Birmingham AI course material.
\end{itemize}

\begin{repl}
: help matches
: ref regexp
: teach json
\end{repl}

The same corpus is published as a website by a Pop-11 program in the tree
(\code{tools/gen-docs.sh}) at \code{https://iotone.github.io/poplog/}, with an
\code{llms.txt} for AI assistants.

\section{Editing and long-lived sessions}

Because the interesting object is a \emph{running heap}, the tooling in this
fork is built to talk to one rather than to a file.

\begin{itemize}
\item \textbf{VED}, the built-in editor: \code{ved myfile.p}. \code{ENTER l1}
  compiles the current line into the live system, \code{ENTER lmr} the marked
  range, \code{ENTER lcp} the current procedure.
\item \textbf{LSP} (\code{tools/pop11-lsp}): diagnostics come from the real
  compiler with \code{pop\_syntax\_only} set, so your file is type-checked by
  the thing that would compile it, and nothing in it runs. A Neovim plugin
  ships in \code{editors/nvim/}.
\item \textbf{Swank} (\code{tools/pop11-swank}): the other half --- a live
  session an editor can attach to. Output streams as it is produced, mishaps
  arrive as data with real frames, and a runaway loop can be interrupted. The
  Emacs package in \code{editors/emacs/} is the client;
  \code{teach swank} is the walkthrough.
\item \textbf{MCP} (\code{tools/pop11-mcp}): the same idea for agents, and
  the way the general of the Robot Army actually gives orders. A persistent,
  natively compiled session exposed as four tools --- \code{pop11\_eval},
  \code{pop11\_help}, \code{pop11\_checkpoint}, \code{pop11\_state}.
  Register it with \code{claude mcp add --scope project pop11 --
  tools/pop11-mcp} and the agent has the command post in its hands.
\end{itemize}

All three servers sit on one transport, \code{pop/lib/lib/jsonrpc.p}, itself
written in Pop-11.

\begin{callout}{The working pattern}
Start one session and keep it. Define helpers once; call them thereafter.
Redefinition is instant and applies to every later call, so a procedure is
something you \emph{sculpt} in a live image rather than edit-compile-run.
Checkpoint the image when the session has earned something worth keeping.
\end{callout}
